\noindent \inouter{\notagramatical{οἰχομένων}{οἴχομαι}{partir}{part., pres., voz méd., gen. pl. masc.}}\numeroverso{1}Τάδε μὲν Περσῶν τῶν οἰχομένων\par
\noindent \numeroverso{}Ἑλλάδ᾽ ἐς αἶαν πιστὰ καλεῖται,\par
\noindent \numeroverso{}καὶ τῶν ἀφνεῶν καὶ πολυχρύσων\par
\noindent \numeroverso{}ἑδράνων φύλακες, κατὰ πρεσβείαν\par
\noindent \numeroverso{5}οὓς αὐτὸς ἄναξ Ξέρξης βασιλεὺς\par
\noindent \numeroverso{}Δαρειογενὴς\par
\noindent \numeroverso{}εἵλετο χώρας ἐφορεύειν.\par
\noindent \inouter{\notagramatical{νόστῳ}{νόστος/νόστου}{retorno}{subst., dat. sing. masc.}}\numeroverso{}ἀμφὶ δὲ νόστῳ τῷ βασιλείῳ\par
\noindent \numeroverso{}καὶ πολυχρύσου στρατιᾶς ἤδη\par
\noindent \numeroverso{10}κακόμαντις ἄγαν ὀρσολοπεῖται\par
\noindent \numeroverso{}θυμὸς ἔσωθεν.\par
\noindent \numeroverso{}πᾶσα γὰρ ἰσχὺς Ἀσιατογενὴς\par
\noindent \numeroverso{}ᾤχωκε, νέον δ᾽ ἄνδρα βαΰζει,\par
\noindent \numeroverso{}κοὔτε τις ἄγγελος οὔτε τις ἱππεὺς\par
\noindent \numeroverso{15}ἄστυ τὸ Περσῶν ἀφικνεῖται:\par
\noindent \numeroverso{}οἵτε τὸ Σούσων ἠδ᾽ Ἀγβατάνων\par
\noindent \numeroverso{}καὶ τὸ παλαιὸν Κίσσιον ἕρκος\par
\noindent \numeroverso{}προλιπόντες ἔβαν, τοὶ μὲν ἐφ᾽ ἵππων,\par
\noindent \numeroverso{}τοὶ δ᾽ ἐπὶ ναῶν, πεζοί τε βάδην\par
\noindent \numeroverso{20}πολέμου στῖφος παρέχοντες:\par
\noindent \numeroverso{}οἷος Ἀμίστρης ἠδ᾽ Ἀρταφρένης\par
\noindent \numeroverso{}καὶ Μεγαβάτης ἠδ᾽ Ἀστάσπης,\par
\noindent \numeroverso{}ταγοὶ Περσῶν,\par
\noindent \numeroverso{}βασιλῆς βασιλέως ὕποχοι μεγάλου,\par
\noindent \numeroverso{25}σοῦνται, στρατιᾶς πολλῆς ἔφοροι,\par
\noindent \inouter{\notagramatical{ἠδ᾽}{ἠδέ}{e}{conj.}}\numeroverso{}τοξοδάμαντές τ᾽ ἠδ᾽ ἱπποβάται,\par
\noindent \numeroverso{}φοβεροὶ μὲν ἰδεῖν, δεινοὶ δὲ μάχην\par
\noindent \numeroverso{}ψυχῆς εὐτλήμονι δόξῃ:\par
\noindent \numeroverso{}Ἀρτεμβάρης θ᾽ ἱππιοχάρμης\par
\noindent \numeroverso{30}καὶ Μασίστης, ὅ τε τοξοδάμας\par
